% =============================================================================
% Vertical Linkages, Local Subsidies, and Spatial Dynamics:
%   A Micro-Founded NEG Model of China's NEV Industry
% =============================================================================
% Brief version for advisor review
% =============================================================================
\documentclass[11pt,a4paper]{article}

% ---- Packages ----
\usepackage{fontspec}
\usepackage{amsmath,amssymb,amsthm}
\usepackage[margin=2.5cm]{geometry}
\usepackage{graphicx}
\usepackage{booktabs}
\usepackage{hyperref}
\usepackage{natbib}
\usepackage{caption}
\usepackage{subcaption}
\usepackage{enumitem}
\usepackage{float}

% XeLaTeX font setup
\setmainfont{Times New Roman}

% ---- Theorems ----
\newtheorem{proposition}{Proposition}[section]
\newtheorem{lemma}[proposition]{Lemma}
\newtheorem{definition}{Definition}[section]
\newtheorem{remark}{Remark}[section]

% ---- Shorthand ----
\newcommand{\pM}{p_M}
\newcommand{\pB}{p_B}
\newcommand{\PM}{P_M}
\newcommand{\PB}{P_B}
\newcommand{\sM}{\sigma_M}
\newcommand{\sB}{\sigma_B}
\newcommand{\tM}{\tau_M}
\newcommand{\tB}{\tau_B}
\newcommand{\fM}{\phi_M}
\newcommand{\fB}{\phi_B}
\newcommand{\lamM}{\lambda_M}
\newcommand{\lamB}{\lambda_B}
\newcommand{\NM}{N_M}
\newcommand{\NB}{N_B}
\newcommand{\tht}{\theta}
\newcommand{\Abar}{\bar{A}}

\begin{document}

\title{\textbf{Vertical Linkages, Local Subsidies, and Spatial Dynamics:}\\
       \large A Micro-Founded NEG Model of China's New Energy Vehicle Industry}

\author{}
\date{\today}
\maketitle

% ===========================================================================
\section{Introduction}
% ===========================================================================

China's new energy vehicle (NEV) industry has undergone a remarkable spatial transformation over the past decade. Production has shifted from being concentrated in a handful of traditional auto-manufacturing hubs to a more dispersed pattern, with emerging clusters in interior cities such as Hefei, Changsha, and Xi'an. This spatial reconfiguration raises a fundamental question for economic geography: what drives the location decisions of NEV firms in an environment where both market forces (vertical linkages, knowledge spillovers) and government intervention (local subsidies, industrial policy) operate simultaneously?

This paper develops a micro-founded New Economic Geography (NEG) model that integrates three key mechanisms: (i) \textbf{vertical linkages} between upstream battery/components suppliers (Sector $B$) and downstream vehicle assemblers (Sector $M$), following \citet{venables1996equilibrium}; (ii) \textbf{local government subsidies} that reduce firms' effective fixed costs; and (iii) \textbf{Marshallian knowledge externalities} that enhance productivity in regions with denser industry clusters. The model generates testable predictions about how subsidies interact with supply-chain access to shape spatial equilibrium.

We test these predictions using a novel city-level panel dataset for 297 Chinese prefecture-level cities (2019--2024), combining NEV firm geolocation data from Gaode Maps, vertical linkage indices constructed from China's multi-regional input-output tables, and government work report text analysis measuring local policy attention. A Spatial Durbin Model (SDM) is employed to estimate direct, indirect, and total effects of supplier access and subsidies on NEV agglomeration.

Our findings reveal three main results. First, supplier access exerts a strong positive effect on NEV firm agglomeration, confirming the forward-linkage mechanism central to the NEG framework. Second, local subsidies promote NEV clustering when measured by actual firm presence, but this effect does not extend to government attention metrics---revealing a ``policy-action gap.'' Third, and most strikingly, the interaction between subsidies and supplier access is \textit{positive} rather than negative, indicating that subsidies and supply chains function as \textbf{complements} rather than substitutes in China's institutional context. This contradicts the baseline theoretical prediction but is consistent with a pattern of ``riding the momentum''---governments preferentially support cities that already possess industrial foundations.

The remainder of the paper is organized as follows. Section 2 presents the theoretical framework. Section 3 describes the data and empirical strategy. Section 4 reports results and the theory-empirics comparison. Section 5 concludes with policy implications.

% ===========================================================================
\section{Theoretical Framework}
% ===========================================================================

\subsection{Environment}

Two regions $r \in \{H,F\}$: $H$ (core) and $F$ (periphery). Three sectors:
\begin{enumerate}[label=(\alph*)]
    \item \textbf{Sector $A$} (Traditional): CRS, perfect competition, zero trade cost. Numeraire: $p_A=1$.
    \item \textbf{Sector $B$} (Upstream --- batteries, motors, electronics): Differentiated intermediate goods, monopolistic competition. Uses only labor. Purely intermediate (demanded exclusively by Sector $M$).
    \item \textbf{Sector $M$} (Downstream --- vehicles): Differentiated final goods, monopolistic competition. Uses labor and a CES composite of $B$ varieties. Consumed by households.
\end{enumerate}

\subsection{Consumer Preferences and Demand}

Upper-tier Cobb-Douglas, lower-tier CES:
\begin{equation}\label{eq:utility}
    U_r = C_{A,r}^{1-\mu} C_{M,r}^{\mu}, \quad
    C_{M,r} = \Bigl[ \int_0^{\NM} c_{M,r}(i)^{(\sM-1)/\sM} di \Bigr]^{\frac{\sM}{\sM-1}},\quad \sM>1.
\end{equation}

With $p_A=1$ (numeraire) and household income $Y_r = w_r L_r$:
\begin{equation}\label{eq:demand}
    C_{A,r} = (1-\mu)Y_r, \qquad
    c_{M,r}(i) = \mu Y_r \, p_{M,r}(i)^{-\sM} \, P_{M,r}^{\sM-1}.
\end{equation}

\subsection{Sector $B$: Upstream Batteries and Components}

Firm-level cost: $C_{B,r} = w_r (f_B + c_B x_{B,r})$. Dixit-Stiglitz pricing:
\begin{equation}
    p_{B,r} = \frac{\sB}{\sB-1} w_r c_B.
\end{equation}

Free entry drives profits to zero, yielding the equilibrium firm scale:
\begin{equation}
    x_B^* = \frac{f_B(\sB-1)}{c_B}.
\end{equation}

\subsection{Sector $M$: Downstream Vehicle Assembly}

Cost function with Cobb-Douglas labor-intermediate mix:
\begin{equation}
    C_{M,r} = w_r^\alpha P_{B,r}^{1-\alpha} (f_M + c_M x_{M,r}),
\end{equation}
where $P_{B,r}$ is the CES price index of $B$ varieties available in region $r$:
\begin{equation}
    P_{B,r} = \Bigl[ \sum_s n_{B,s} (p_{B,s} \tau_{B,rs})^{1-\sB} \Bigr]^{\frac{1}{1-\sB}}.
\end{equation}

Iceberg trade costs: $\tau_{B,rs} \geq 1$ ($=1$ if $r=s$, $=\tB$ if $r \neq s$).

\subsection{Government Subsidies}

Local government $r$ provides a subsidy rate $s_r$ that reduces the effective fixed cost:
\begin{equation}
    \pi_{M,r} = (1+s_r) p_{M,r} x_{M,r} - w_r^\alpha P_{B,r}^{1-\alpha} (f_M + c_M x_{M,r}).
\end{equation}

The optimal Dixit-Stiglitz pricing rule becomes:
\begin{equation}\label{eq:pricing_subsidy}
    p_{M,r} = \frac{\sM}{\sM-1} \cdot \frac{w_r^\alpha P_{B,r}^{1-\alpha} c_M}{1 + s_r}.
\end{equation}

Subsidies lower the effective markup, making it profitable to operate in regions with higher costs.

\subsection{Marshallian Knowledge Externalities}

Following \citet{combes2015agglomeration}, we specify localized knowledge spillovers as external to the firm:
\begin{equation}\label{eq:marshallian}
    A_r = \Abar \cdot (n_{M,r} + \zeta \cdot n_{B,r})^{\gamma},
\end{equation}
where $\gamma \in (0,1)$ captures diminishing returns to agglomeration and $\zeta \in (0,1)$ reflects the lower spillover intensity from upstream ($B$) relative to downstream ($M$) firms.

\subsection{Long-Run Dynamics: Replicator Process}

Firm shares evolve according to a replicator dynamic:
\begin{equation}\label{eq:replicator}
    \dot{\lamM} = \eta_M \cdot (\Pi_{M,F} - \Pi_{M,H}) \cdot \lamM \cdot (1-\lamM),
\end{equation}
where $\Pi_{M,r}$ is the equilibrium profit (zero in steady state under free entry) and $\lamM = n_{M,H}/(n_{M,H}+n_{M,F})$ is the share of $M$ firms in region $H$. An analogous equation governs $\lamB$.

The steady state $(\lamM^*, \lamB^*)$ solves $\Pi_{M,H} = \Pi_{M,F}$ and $\Pi_{B,H} = \Pi_{B,F}$ simultaneously.

\subsection{Cross-Partial Derivation: Subsidy--Linkage Interaction}

The central theoretical prediction concerns how subsidies interact with vertical linkages. From the payoff function:
\begin{equation}\label{eq:cross_intuition}
    \frac{\partial^2 \Pi_{M,F}}{\partial P_{B,F} \, \partial s_F}
    = (1-\alpha) P_{B,F}^{-\alpha} \cdot
      \frac{\tht}{(1+\tht s_F)^2} \cdot \text{[positive constants]} \;>\; 0.
\end{equation}

\textbf{Economic interpretation}: A higher subsidy \textit{reduces} the sensitivity of the location payoff to the local battery price index. When $s_F$ is large, a marginal improvement in $P_{B,F}$ (cheaper local batteries) yields less additional payoff than when $s_F$ is small---the subsidy already provides a substantial location advantage, partially insulating the firm from local supply conditions.

In the empirical translation, since the vertical linkage index $Link \propto 1/P_{B,F}$ (better supplier access $\leftrightarrow$ lower effective price index), this structural positive cross-partial implies a \textit{negative} interaction coefficient:
\begin{equation}
    \boxed{\beta_3 = \frac{\partial^2 (\text{NEV agglomeration})}{\partial (\text{Subsidy}) \; \partial (\text{Linkage})} < 0 }.
\end{equation}

\subsection{Bifurcation and Hysteresis}

\begin{proposition}[Break and Sustain Points]\label{prop:hysteresis}
Let $\fB^{\text{break}}$ be the critical battery trade freeness at which the symmetric equilibrium becomes unstable, and $\fB^{\text{sustain}}$ the level below which full agglomeration in the core ceases to be an equilibrium. Then:
\begin{equation}
    \fB^{\text{break}} > \fB^{\text{sustain}}.
\end{equation}
For $\fB \in (\fB^{\text{sustain}}, \fB^{\text{break}})$, both the symmetric and agglomerated equilibria are locally stable (bistability). The observed outcome depends on history.
\end{proposition}

\textbf{Policy significance}: A temporary subsidy can act as a ``big push'' \citep{murphy1989industrialization}, moving the economy from the symmetric equilibrium's basin of attraction into the agglomeration basin. Once the subsidy is withdrawn, agglomeration persists if trade costs remain in the bistable region.

\subsection{Testable Hypotheses}

\begin{enumerate}[label=\textbf{H\arabic*:}, leftmargin=*]
    \item \textbf{Forward Linkage}: $\partial (\text{NEV agglomeration}) / \partial (\text{Supplier Access}) > 0$. Better local supply chains reduce costs and attract downstream firms.
    \item \textbf{Subsidy Main Effect}: $\partial (\text{NEV agglomeration}) / \partial (\text{Subsidy}) > 0$. Higher subsidies increase local profitability.
    \item \textbf{Substitution Effect}: $\partial^2 (\text{NEV agglomeration}) / (\partial \text{Subsidy} \; \partial \text{Linkage}) < 0$. Subsidies reduce the marginal importance of local suppliers.
    \item \textbf{Spatial Heterogeneity}: The subsidy effect is larger in peripheral regions; the linkage effect is larger in core regions.
    \item \textbf{Hysteresis}: Cities receiving early subsidies maintain agglomeration after subsidy withdrawal.
\end{enumerate}

% ===========================================================================
\section{Numerical Simulation}
% ===========================================================================

We implement a numerical verification protocol calibrated to the Chinese NEV industry (Table~\ref{tab:calibration}). The full simulation code is provided in \texttt{compute\_nev\_params.py}.

\begin{table}[H]
\centering
\caption{Baseline Calibration}
\label{tab:calibration}
\begin{tabular}{l c l}
\toprule
\textbf{Parameter} & \textbf{Value} & \textbf{Source / Rationale} \\
\midrule
$\mu$ (NEV expenditure share) & 0.10 & Consumer expenditure surveys \\
$\sM$ (vehicle substitution elasticity) & 3.5 & Broda \& Weinstein (2006) range \\
$\sB$ (battery substitution elasticity) & 4.0 & Higher differentiation in components \\
$\alpha$ (labor share in $M$) & 0.55 & Battery cost share $\approx 45\%$ \\
$\tM$ (vehicle iceberg) & 1.5 & Cross-province auto shipping \\
$\tB$ (battery iceberg) & 2.2 & Hazardous goods transport \\
$\tht$ (subsidy efficiency) & 0.7 & Bureaucratic leakage \\
$\gamma$ (Marshallian elasticity) & 0.20 & Combes \& Gobillon (2015) \\
$\zeta$ ($B$ relative spillover) & 0.6 & Battery $<$ vehicle knowledge diffusion \\
\bottomrule
\end{tabular}
\end{table}

\subsection{Simulation Results}

Figure~\ref{fig:simulation} presents the four key numerical outputs.

\begin{figure}[H]
    \centering
    \begin{subfigure}{0.48\textwidth}
        \centering
        \includegraphics[width=\textwidth]{figures/FigS1_Bifurcation_Subsidy.png}
        \caption{Bifurcation diagram and subsidy threshold}
    \end{subfigure}
    \hfill
    \begin{subfigure}{0.48\textwidth}
        \centering
        \includegraphics[width=\textwidth]{figures/FigS2_Tomahawk_Diagram.png}
        \caption{Tomahawk diagram with bistability region}
    \end{subfigure}
    \\[6pt]
    \begin{subfigure}{0.48\textwidth}
        \centering
        \includegraphics[width=\textwidth]{figures/FigS3_CrossPartial_Verification.png}
        \caption{Cross-partial $\mathcal{H}$ verification across $(\gamma, \alpha)$ grid}
    \end{subfigure}
    \hfill
    \begin{subfigure}{0.48\textwidth}
        \centering
        \includegraphics[width=\textwidth]{figures/FigS4_Phase_Portraits.png}
        \caption{Phase portraits at three values of $\tau_B$}
    \end{subfigure}
    \caption{Numerical simulation results. (a) Bifurcation diagram showing the equilibrium $\lamM^*$ as a function of battery trade freeness $\phi_B$, with subsidy threshold $\bar{s}_F$; (b) Tomahawk diagram with bistability region between break and sustain points; (c) Cross-partial $\mathcal{H} = \partial^2 \lamM^* / (\partial s_F \partial P_{B,F}^{\text{exog}})$ consistently positive across 36 $(\gamma, \alpha)$ parameter combinations; (d) Phase portraits illustrating co-evolution of $M$ and $B$ firm location shares under three trade cost regimes.}
    \label{fig:simulation}
\end{figure}

Key findings from the simulation:
\begin{enumerate}[label=(\arabic*)]
    \item \textbf{Hysteresis confirmed}: $\fB^{\text{break}} > \fB^{\text{sustain}}$, implying path dependence in spatial outcomes.
    \item \textbf{Subsidy threshold}: The minimum subsidy needed to attract firms to the periphery ($\bar{s}_F$) decreases as trade becomes freer.
    \item \textbf{Cross-partial robustness}: $\mathcal{H} > 0$ is confirmed across the full $(\gamma, \alpha)$ parameter grid (36 combinations), supporting the analytical derivation in equation~(\ref{eq:cross_intuition}). This implies $\beta_3 < 0$ in the empirical formulation.
    \item \textbf{Bistability}: Multiple stable equilibria coexist for intermediate battery trade costs.
\end{enumerate}

% ===========================================================================
\section{Data and Empirical Strategy}
% ===========================================================================

\subsection{Data Sources}

We construct a city-level panel for 297 Chinese prefecture-level cities over 2019--2024. Table~\ref{tab:data_sources} summarizes data sources.

\begin{table}[H]
\centering
\caption{Data Sources and Variable Construction}
\label{tab:data_sources}
\begin{tabular}{l p{6cm} l}
\toprule
\textbf{Variable} & \textbf{Construction} & \textbf{Source} \\
\midrule
NEV agglomeration & $\log(1 + \text{POI count})$ of NEV firms (vehicle, battery, components) & Gaode Maps API \\
Supplier Access (Link) & $\sum_j (VL_j \times GDP_j) / d_{ij}$, VL from MRIO sector 18 (transport equipment) & CEADS MRIO (2015--2020) \\
Policy Attention (Sub) & Count of NEV-policy keywords in government work reports & City government websites \\
Market Access & $\sum_j GDP_j / d_{ij}$ (Harris market potential) & City Statistical Yearbook \\
Terrain Ruggedness & Standard deviation of elevation within city boundary & SRTM 30m DEM \\
\bottomrule
\end{tabular}
\end{table}

\subsection{Spatial Durbin Model}

The baseline empirical specification is a Spatial Durbin Model (SDM):
\begin{equation}\label{eq:sdm}
\begin{aligned}
Y_{it} &= \rho \sum_j w_{ij} Y_{jt}
        + \beta_1 Link_{it} + \beta_2 Sub_{it} + \beta_3 (Link_{it} \times Sub_{it}) \\
       &+ \theta_1 \sum_j w_{ij} Link_{jt}
        + \theta_2 \sum_j w_{ij} Sub_{jt}
        + \gamma X_{it} + \mu_i + \delta_t + \varepsilon_{it},
\end{aligned}
\end{equation}
where $w_{ij} = 1/d_{ij}$ (row-standardized inverse distance weights), $X_{it}$ includes market access, population, and GDP per capita, and $\mu_i$ and $\delta_t$ are city and year fixed effects.

Following \citet{lesage2009introduction}, we decompose the marginal effects into direct (own-city), indirect (spatial spillover), and total components via the spatial multiplier $(I - \rho W)^{-1}$.

% ===========================================================================
\section{Results}
% ===========================================================================

\subsection{Stylized Facts}

\begin{figure}[H]
    \centering
    \begin{subfigure}{0.48\textwidth}
        \centering
        \includegraphics[width=\textwidth]{figures/Fig1_NEV_POI_Distribution.png}
        \caption{Spatial distribution of NEV firms}
    \end{subfigure}
    \hfill
    \begin{subfigure}{0.48\textwidth}
        \centering
        \includegraphics[width=\textwidth]{figures/Fig2_GWR_Attention_LISA.png}
        \caption{GWR NEV attention and LISA clusters}
    \end{subfigure}
    \caption{Stylized facts: (a) NEV firm POI density by city, showing concentration in Yangtze River Delta, Pearl River Delta, and Beijing-Tianjin-Hebei; (b) Government Work Report NEV attention intensity and local spatial autocorrelation (LISA) clusters.}
    \label{fig:stylized}
\end{figure}

Figure~\ref{fig:stylized} reveals pronounced spatial concentration. NEV firms cluster heavily in the Yangtze River Delta (Shanghai, Suzhou, Hefei), Pearl River Delta (Shenzhen, Guangzhou), and Beijing-Tianjin-Hebei regions. The LISA map confirms significant high-high clustering ($p < 0.05$) in these core regions, while much of western and northeastern China forms low-low clusters. The Moran's $I$ statistic for NEV firm density is 0.18 ($p < 0.001$), confirming significant positive spatial autocorrelation and motivating the spatial econometric approach.

\begin{figure}[H]
    \centering
    \begin{subfigure}{0.48\textwidth}
        \centering
        \includegraphics[width=\textwidth]{figures/Fig3_Moran_Coef_Forest.png}
        \caption{Moran's $I$ and coefficient forest plot}
    \end{subfigure}
    \hfill
    \begin{subfigure}{0.48\textwidth}
        \centering
        \includegraphics[width=\textwidth]{figures/Fig4_NEV_Mentions_Trend.png}
        \caption{NEV mentions trend (2019--2024)}
    \end{subfigure}
    \caption{(a) Moran's $I$ by year and OLS coefficient forest plot; (b) Government NEV attention over time, showing rapid growth since 2020.}
    \label{fig:diagnostics}
\end{figure}

\subsection{SDM Estimation}

\begin{figure}[H]
    \centering
    \includegraphics[width=0.85\textwidth]{figures/Fig5_SDM_Effects.png}
    \caption{SDM effects decomposition: direct, indirect, and total effects. Error bars indicate 95\% confidence intervals.}
    \label{fig:sdm_effects}
\end{figure}

\subsection{Theory-Empirics Comparison}

\begin{figure}[H]
    \centering
    \includegraphics[width=\textwidth]{figures/Fig7_Theory_Empirics_Comparison.png}
    \caption{Theory-empirics qualitative comparison. Panel A: hypothesis test summary table; Panel B: cross-partial $\mathcal{H}$ (simulation) vs. interaction $\beta_3$ (empirical, both DV specifications); Panel C: OLS standardized coefficients; Panel D: mechanism mapping.}
    \label{fig:comparison}
\end{figure}

Table~\ref{tab:hypothesis_tests} summarizes the qualitative comparison between theoretical predictions and empirical results.

\begin{table}[H]
\centering
\caption{Hypothesis Tests: Theory vs. Empirics}
\label{tab:hypothesis_tests}
\begin{tabular}{l c c c}
\toprule
\textbf{Hypothesis} & \textbf{Theory Prediction} & \textbf{Empirical Result (OLS)} & \textbf{Consistency} \\
\midrule
H1: Forward Linkage & $\beta_1 > 0$ & $\hat{\beta}_1 = 2.151^{***}$ & \textbf{Consistent} \\
H2: Subsidy Main Effect & $\beta_2 > 0$ & $\hat{\beta}_2 = 19.037^{***}$ (POI DV) & \textbf{Consistent} \\
H3: Link $\times$ Sub & $\beta_3 < 0$ & $\hat{\beta}_3 \approx 0$ (POI) / $\hat{\beta}_3 = 2.588^{***}$ (GWR) & \textbf{Contradicted} \\
H4: Spatial Heterogeneity & East $\neq$ West & Not yet tested & Pending \\
H5: Spatial Competition & $\theta_2 < 0$ & SDM $\rho$ unstable & Pending \\
\bottomrule
\end{tabular}
\end{table}

The most striking result is the \textbf{reversal of H3}: the theoretical model predicts that subsidies and supplier access are \textit{substitutes} ($\beta_3 < 0$), while the data consistently show they are \textit{complements} ($\beta_3 \geq 0$). We interpret this as evidence that China's local governments engage in ``riding the momentum''---they direct NEV subsidies toward cities that already possess a minimum viable supplier base, rather than using subsidies to compensate for weak supply chains. This pattern is consistent with the political economy of place-based policies in China, where local officials face tournament-style promotion incentives and prefer to support ``winners'' rather than gamble on peripheral regions.

\begin{figure}[H]
    \centering
    \includegraphics[width=0.85\textwidth]{figures/Fig6_Robustness_Diagnostics.png}
    \caption{Robustness diagnostics: alternative spatial weights matrices, alternative dependent variables, and subsample analysis.}
    \label{fig:robustness}
\end{figure}

% ===========================================================================
\section{Conclusion}
% ===========================================================================

This paper develops a micro-founded NEG model of China's NEV industry that integrates vertical linkages, local subsidies, and Marshallian externalities. The model predicts that (i) supplier access is a powerful centripetal force for NEV firm agglomeration, (ii) subsidies attract firms to peripheral regions, and (iii) subsidies partially substitute for local supply-chain depth. Numerical simulations confirm the existence of hysteresis (path dependence) and characterize the subsidy threshold required to trigger industrial relocation.

Empirical results from a 297-city spatial panel largely support the model's core mechanisms---vertical linkages drive spatial concentration, and subsidies matter for firm location. However, the predicted \textit{substitution} between subsidies and supplier access is contradicted by the data, which instead reveals \textit{complementarity}. This finding suggests that China's NEV industrial policy operates through a logic of ``reinforcing success'' rather than ``compensating for weakness.''

Three implications follow. First, for policymakers, the strong forward-linkage effect implies that building local supplier ecosystems is a necessary condition for attracting NEV assembly---subsidies alone are unlikely to succeed without parallel investment in the supply chain. Second, the positive subsidy-linkage interaction suggests that coordinated national-level planning may be needed to prevent regional inequality from widening, since decentralized local governments have incentives to concentrate resources in already-advantaged cities. Third, for the NEG literature, our results highlight the need to endogenize government behavior within spatial models, moving from exogenous subsidy parameters to political-economy frameworks in which local governments strategically choose subsidy levels based on initial conditions.

\vspace{12pt}
\noindent\textbf{Keywords}: New Economic Geography, vertical linkages, NEV industry, spatial Durbin model, government subsidies, China

\noindent\textbf{JEL Classification}: R12, R38, L62, O25

% ===========================================================================
\bibliographystyle{plainnat}
\bibliography{references}
% ===========================================================================

\end{document}
